% Part: first-order-logic
% Chapter: natural-deduction
% Section: rules-and-proofs

\documentclass[../../../include/open-logic-section]{subfiles}

\begin{document}

\iftag{FOL}
      {\olfileid{fol}{ntd}{rul}}
      {\olfileid{pl}{ntd}{rul}}

\olsection{Regels en \usetoken{P}{derivation}}

\begin{explain}
Een systeem van natuurlijke deductie volgt zo dicht mogelijk de
informele manier van redeneren in een wiskundig bewijs. Daarom heet
het ``natuurlijk''. Een bewijs begint met aannamen. Daarna passen we
afleidingsregels toe. De regels \Intro{\lnot}, \Intro{\lif},
\iftag{FOL}{\Elim{\lor} en \Elim{\lexists}}{ en \Elim{\lor}} trekken
aannamen in: zulke aannamen worden ``!!{discharged}''. Naast de
afleidingsstap staat het label van de !!{discharged} aanname, zodat
duidelijk is welke aanname die stap intrekt.
\end{explain}

\begin{defn}[Aanname]
Een \emph{aanname} is elke !!{sentence} die bovenaan, dus als eerste
knoop, in een tak staat.
\end{defn}

!!^{derivation}s in natuurlijke deductie zijn bomen van
!!{sentence}s. De bovenste !!{sentence}s zijn aannamen. Staat
!!a{sentence} onder
één, twee of drie sequenten, dan moet een afleidingsregel hem correct
uit die sequenten laten volgen. De !!{sentence}s boven de regel heten
de \emph{premissen}. De !!{sentence} eronder heet de
\emph{conclusie}. Voor elke !!{operator} vormen twee regels een paar:
een introductieregel en een eliminatieregel. De eerste voert
!!a{operator} in de conclusie in; de tweede verwijdert !!a{operator}
uit een premisse. Sommige regels staan toe dat een aanname van een
bepaald type wordt \emph{!!{discharged}}. We geven de aanname en de
afleidingsstap dan hetzelfde label. Zo is zichtbaar welke
afleidingsstap welke aanname wordt !!{discharged}. De gelabelde aanname
schrijven we als ``$\Discharge{!A}{n}$.''

% Only include this sentence is on of \land, lor, lif or lnot is defined.

\iftag{notprvNot,notprvAnd,notprvOr,notprvIf}{}%
	{Gewoonlijk geven we regels voor alle !!{operator}s $\land$, $\lor$, $\lif$, $\lnot$ en $\lfalse$, ook als we sommige daarvan al via andere operatoren hebben gedefinieerd.}
\end{document}
